\SourceLeaf{31} kleinen Zeitabschnitt $\mathrm{d}t$ einnehmen würden, wenn sie alle vollkommen frei wären, und zwar infolge der während dieser Zeit auf sie wirkenden Kräfte und der zur Zeit $t$ erreichten Geschwindigkeiten und Richtungen. Die tatsächlichen Orte $c,c',c''$ usw. sind dann unter allen mit den Bedingungen des Systems vereinbaren Orten diejenigen, für die $m(bc)^2+m'(b'c')^2+m''(b''c'')^2$ usw. ein Minimum wird.

Das Gleichgewicht ist offenbar nur ein einzelner Fall des allgemeinen Gesetzes. Die Bedingung dafür ist, dass
\[
m(ab)^2+m'(a'b')^2+m''(a''b'')^2\ \text{u. s. w.}
\]
selbst ein Minimum ist oder dass das Verbleiben des Systems im Ruhezustand der freien Bewegung der einzelnen Punkte näher kommt als jedes mögliche Verlassen dieses Zustands.

Unser Prinzip lässt sich aus den oben genannten Prinzipien leicht auf folgende Weise ableiten.

Die auf den materiellen Punkt $m$ wirkende Kraft setzt sich offenbar aus zwei Kräften zusammen. Die erste führt ihn zusammen mit der zur Zeit $t$ bestehenden Geschwindigkeit und Richtung in der Zeit $\mathrm{d}t$ von $a$ nach $c$. Die zweite würde ihn in derselben Zeit aus der Ruhe in $c$ entlang $cb$ führen, wenn man den Punkt als frei betrachtet. Dasselbe gilt für die anderen Punkte. Nach dem Prinzip von \textsc{d’Alembert} müssen demnach die Punkte $m,m',m''$ usw. allein unter der Wirkung der zweiten, entlang $cb,c'b',c''b''$ usw. gerichteten Kräfte an den Orten $c,c',c''$ usw. aufgrund der Bedingungen des Systems im Gleichgewicht sein.

\expandafter\def\csname HistoricalNoteText3\endcsname{Die übliche Formulierung setzt stillschweigend Bedingungen voraus, unter denen zu jeder möglichen Bewegung auch die entgegengesetzte möglich ist. Ein Punkt kann zum Beispiel gezwungen sein, auf einer bestimmten Fläche zu bleiben, oder der Abstand zwischen zwei Punkten kann unveränderlich sein und dergleichen. Diese Einschränkung ist aber unnötig und der Natur nicht immer angemessen. Die Oberfläche eines undurchdringlichen Körpers zwingt einen auf ihr befindlichen materiellen Punkt nicht, dort zu bleiben. Sie verhindert lediglich, dass er auf die eine Seite übertritt. Ein gespannter, nicht dehnbarer, aber biegsamer Faden zwischen zwei Punkten macht nur eine Vergrößerung ihres Abstands unmöglich und lässt eine Verkleinerung zu, usw. Warum sollten wir also das Gesetz der virtuellen Geschwindigkeiten nicht lieber gleich zu Beginn so formulieren, dass es \textit{alle} Fälle umfasst?}
Nach dem Prinzip der virtuellen Geschwindigkeiten erfordert dieses Gleichgewicht, dass die Summe der Produkte aus je drei Faktoren, nämlich aus den Massen $m,m',m''$ usw., den Strecken $cb,c'b'\ c''b''$ usw. und beliebigen mit den Bedingungen des Systems vereinbaren Bewegungen dieser Punkte, jeweils auf die entsprechende Strecke projiziert, immer $=0$ ist, wie man es gewöhnlich ausdrückt\HistoricalNote{3}
